\section{Выводы}

В результате выполнения лабораторной работы был успешно спроектирован и реализован словарь на основе АВЛ-дерева. 

\textbf{Об алгоритме.} Практическая реализация балансирующихся деревьев требует аккуратности в пересчете высот и отслеживании указателей при вращениях. Использование прямого обхода дерева для сериализации доказало свою эффективность: чтение бинарного дампа и восстановление структуры занимает линейное время $O(N)$ и исключает необходимость повторной балансировки при запуске программы.

\textbf{О современных стандартах C++.} Использование \texttt{std::unique\_ptr} сделало код намного безопаснее. При возникновении исключительных ситуаций (например, при ошибках формата во время загрузки файла) вся выделенная память освобождается автоматически. Каскадное удаление узлов в деструкторе дерева избавляет от ручного управления памятью.

\textbf{О производительности.} Бенчмарк выявил интересную особенность: несмотря на строгую логарифмическую асимптотику, реальное время работы сильно зависит от системных вызовов. Текущая рекурсивная архитектура с умными указателями накладывает ощутимый оверхед. Это открывает пространство для дальнейших исследований: в будущем целесообразно провести профилирование данной программы, чтобы выявить узкие места инфраструктуры языка и оптимизировать их путем перехода к итеративным обходам и пулам памяти.
\pagebreak